\chapter{Fixed-income foundations}\label{ch:foundations}
\label{sec:present-value}\label{sec:compounding}
Money available today can earn interest, so payment dates affect economic value. Present value expresses a future payment in today's units; its discount factor is today's value of one future unit. Discrete compounding credits interest periodically; a nominal annual rate is divided by the compounding frequency. Continuous compounding uses exponential accumulation.

\label{sec:contract}\label{sec:dates}\label{sec:price-quotation}
A fixed-rate bond is a debt contract paying fixed interest coupons and returning principal at maturity. Face, or par, value is the principal reference amount; redemption repays it. The annual coupon rate determines annual interest, divided into coupon amounts according to payment frequency. Settlement is the valuation and transfer date; maturity is the final payment date. Accrued interest allocates the current coupon to elapsed time. Dirty price includes this accrual; clean price subtracts it \parencite{tuckman2011}.

\section{Yield to maturity versus zero/spot rates versus par yields}\label{sec:zero-par}\label{sec:ytm-concept}\label{sec:term-structure}
Yield to maturity (YTM) is the single rate equating discounted promised payments to dirty price under specified compounding. A zero-coupon bond pays only at redemption: its spot, or zero-coupon, rate applies to that single horizon. For maturity \(T>0\) years and continuous annual decimal zero rate \(r(T)\), the dimensionless discount factor is exactly
\begin{equation}\label{eq:discount}
DF(T)=\exp[-r(T)T].
\end{equation}
A par yield is the annual coupon rate making a specified bond worth face value under given discount factors, conventionally at a coupon-date settlement with no accrual. It depends on the whole payment schedule, whereas a spot rate discounts one maturity. The term structure describes rates across maturities, or tenors; a yield curve represents it under a stated rate convention.

Treasury constant-maturity (CMT)/par-style yields differ from zero rates. Fitting them does not bootstrap zeros: bootstrapping infers discount factors from instrument prices and cash flows \parencite{bis2005}. Here the invented observations have no intrinsic compounding convention; risk deliberately interprets fitted rates as continuous zeros, an educational assumption.

\label{sec:units}
All rates are decimal: \(0.05=5\%\). One basis point (bp) equals \(0.01\) percentage point, or \(0.0001\) decimal.

\section{Literature and methodological position}
The literature separates curve representation from estimation using instruments. Nelson--Siegel provides a compact curve family; Svensson adds flexibility through a second curvature component \parencite{nelson1987,svensson1994}. This flexibility is useful for representing shapes, but it does not establish that six parameters are necessary for these ten illustrative observations. Institutional approaches also differ: Anderson and Sleath use a spline-based method, while G\"urkaynak, Sack and Wright estimate a discount function from Treasury instruments \parencite{anderson2001,gurkaynak2007}. Such studies address instrument selection and estimation choices absent here. The present study assesses an existing NSS implementation; it neither compares competing curve families nor establishes empirical superiority. Its research contribution is therefore reproducible verification under specified conventions.
